\paragrafos{Κύκλος και στοιχεία του}{Β.3.12, «Κύκλος και στοιχεία του»}

\begin{anagke}
Ο μπαμπάς της Ρήνης τοποθετεί έναν δρομολογητή Wi-Fi στο κέντρο της αυλής του σχολείου, και το σήμα του φτάνει ακριβώς $15$ μέτρα προς κάθε κατεύθυνση. Ο Ηλίας ρωτά: «Πού μπορούμε να καθίσουμε και να έχουμε σήμα;». Ο Νίκος τρέχει $15$ μέτρα από τον δρομολογητή και σταματά. Ο Στέλιος πηγαίνει $20$ μέτρα μακριά, και η Φανή μένει $10$ μέτρα μακριά. «Όλοι όσοι είναι στα $15$ μέτρα σχηματίζουν μια γραμμή», λέει η Αγγελική και με ένα σχοινί την ζωγραφίζει στο χώμα: στρογγυλή. Η Ζωή ρωτά πώς λέγονται τα διάφορα κομμάτια του σχήματος και πόσο «μεγάλο» είναι ένα τμήμα του.

Σε αυτή την παράγραφο μαθαίνουμε τον κύκλο και τα στοιχεία του, μετράμε τόξα με μοίρες και φτιάχνουμε σχήματα με τον διαβήτη.
\end{anagke}

\begin{lexi}{κύκλος, ακτίνα, χορδή, τόξο}{τα στοιχεία του κύκλου}
Ο \lexiword{κύκλος} είναι αυτό που «κυλάει». Η \lexiword{ακτίνα} είναι η «αχτίδα», μια γραμμή που ξεκινά από το κέντρο όπως οι αχτίδες του ήλιου. Το \lexiword{τόξο} είναι το όπλο που τεντώνεται, και η \lexiword{χορδή} είναι το νήμα του τόξου, που ενώνει τα δύο άκρα του. Η \lexiword{διάμετρος} είναι αυτό που μετρά «δια μέσου» του κέντρου.
\end{lexi}

\begin{metafrasi}
Όλα τα σημεία σε ίδια απόσταση από το κέντρο. & κύκλος $(\mathrm{Ο},\rho)$\\
Η απόσταση από το κέντρο. & ακτίνα $\rho$\\
Τμήμα που ενώνει δύο σημεία του κύκλου. & χορδή\\
Χορδή που περνά από το κέντρο. & διάμετρος $\delta=2\rho$\\
Μέρος του κύκλου ανάμεσα σε δύο σημεία. & τόξο $\overset{\frown}{\mathrm{ΑΒ}}$\\
\end{metafrasi}

\tmima{Η θεωρία}

\begin{orismos}{κύκλος}
\textbf{Κύκλος} λέγεται το σύνολο όλων των σημείων του επιπέδου που απέχουν την ίδια απόσταση από ένα σταθερό σημείο $\mathrm{Ο}$. Το $\mathrm{Ο}$ λέγεται \textbf{κέντρο} και η κοινή απόσταση \textbf{ακτίνα} $\rho$. Ο κύκλος με κέντρο $\mathrm{Ο}$ και ακτίνα $\rho$ συμβολίζεται $(\mathrm{Ο},\rho)$ και σχεδιάζεται με τον \textbf{διαβήτη}.
\end{orismos}

Δύο κύκλοι είναι \textbf{ίσοι} όταν έχουν ίσες ακτίνες. Δύο κύκλοι με το ίδιο κέντρο λέγονται \textbf{ομόκεντροι}.

\begin{center}
\begin{tikzpicture}[scale=0.95]
\draw[thick] (0,0) circle (1.8);
\fill (0,0) circle (1.8pt) node[below left] {\small $\mathrm{Ο}$};
\draw[gblue,thick] (0,0)--(1.27,1.27) node[midway,above left] {\small $\rho$};
\fill (1.27,1.27) circle (1.6pt) node[above right] {$\mathrm{Α}$};
\draw (-1.8,0)--(1.8,0); \fill (-1.8,0) circle (1.6pt) node[left] {$\mathrm{Β}$}; \fill (1.8,0) circle (1.6pt) node[right] {$\mathrm{Γ}$};
\node[above] at (-0.7,0.02) {\small διάμετρος};
\draw[orange!80!black,thick] (-1.5,-0.995)--(0.9,-1.56);
\fill (-1.5,-0.995) circle (1.6pt) node[left] {$\mathrm{Δ}$}; \fill (0.9,-1.56) circle (1.6pt) node[below right] {$\mathrm{Ε}$};
\node[orange!80!black] at (-0.2,-1.0) {\small χορδή};
\draw[red!70!black,very thick] (1.8,0) arc (0:45:1.8);
\node[red!70!black] at (2.35,0.75) {\small τόξο};
\end{tikzpicture}
\end{center}

\begin{itemize}[leftmargin=1.2em,itemsep=0.3ex,topsep=2pt]
\item \textbf{Χορδή} είναι το ευθύγραμμο τμήμα που ενώνει δύο σημεία του κύκλου.
\item \textbf{Διάμετρος} είναι η χορδή που περνά από το κέντρο. Είναι η μεγαλύτερη χορδή, ίση με το διπλάσιο της ακτίνας ($\delta=2\rho$), και το κέντρο είναι το μέσο της.
\item Δύο σημεία $\mathrm{Α}$, $\mathrm{Β}$ του κύκλου τον χωρίζουν σε δύο μέρη, καθένα από τα οποία λέγεται \textbf{τόξο} με άκρα τα $\mathrm{Α}$ και $\mathrm{Β}$, $\overset{\frown}{\mathrm{ΑΒ}}$. Ένα τόξο που είναι το μισό του κύκλου λέγεται \textbf{ημικύκλιο}.
\item Ο \textbf{κυκλικός δίσκος} είναι ο κύκλος μαζί με το εσωτερικό του. Ένα σημείο $\mathrm{Μ}$ του δίσκου έχει $\mathrm{ΟΜ}\le\rho$: πάνω στον κύκλο όταν $\mathrm{ΟΜ}=\rho$, μέσα του όταν $\mathrm{ΟΜ}<\rho$, έξω όταν $\mathrm{ΟΜ}>\rho$.
\end{itemize}

\begin{orismos}{επίκεντρη γωνία και μέτρο τόξου}
Μια γωνία λέγεται \textbf{επίκεντρη} όταν η κορυφή της είναι το κέντρο του κύκλου. Οι πλευρές της τέμνουν τον κύκλο στα $\mathrm{Α}$, $\mathrm{Β}$ και το τόξο $\overset{\frown}{\mathrm{ΑΒ}}$ που περιέχεται στο εσωτερικό της λέγεται \textbf{αντίστοιχο τόξο}. Το \textbf{μέτρο ενός τόξου} (σε μοίρες) είναι ίσο με το μέτρο της αντίστοιχης επίκεντρης γωνίας.
\end{orismos}

\begin{center}
\begin{tikzpicture}[scale=0.95]
\draw[thick] (0,0) circle (1.7);
\fill (0,0) circle (1.8pt) node[below] {\small $\mathrm{Ο}$};
\draw[thick] (0,0)--(1.7,0) (0,0)--(0.85,1.47);
\fill (1.7,0) circle (1.6pt) node[right] {$\mathrm{Α}$}; \fill (0.85,1.47) circle (1.6pt) node[above right] {$\mathrm{Β}$};
\draw[red!70!black,very thick] (1.7,0) arc (0:60:1.7);
\draw[gblue] (0.45,0) arc (0:60:0.45); \node[gblue] at (0.72,0.3) {\small $\omega$};
\end{tikzpicture}
\end{center}

\begin{idiotita}{τόξα και επίκεντρες γωνίες}
\begin{itemize}[leftmargin=1.2em,itemsep=0.3ex,topsep=2pt]
\item Ο κύκλος είναι ένα τόξο $360^\circ$ και το ημικύκλιο ένα τόξο $180^\circ$.
\item Δύο κάθετες διάμετροι χωρίζουν τον κύκλο σε τέσσερα ίσα τόξα των $90^\circ$.
\item Στον ίδιο κύκλο (ή σε ίσους κύκλους) δύο ίσες επίκεντρες γωνίες έχουν ίσα αντίστοιχα τόξα και αντίστροφα: δύο ίσα τόξα έχουν ίσες επίκεντρες γωνίες.
\end{itemize}
\end{idiotita}

Ο διαβήτης είναι και εργαλείο κατασκευής. Αν θέλουμε τρίγωνο με πλευρές $a,b,c$, σχεδιάζουμε ένα τμήμα $\mathrm{ΑΒ}=c$ και δύο κύκλους, $(\mathrm{Α},b)$ και $(\mathrm{Β},a)$. Το σημείο τομής τους είναι η τρίτη κορυφή.

\begin{prosoxi}
\begin{itemize}[leftmargin=1.2em,itemsep=0.3ex,topsep=2pt]
\item Ο κύκλος είναι μόνο η γραμμή· το εσωτερικό μαζί με τη γραμμή είναι ο κυκλικός δίσκος.
\item Κάθε διάμετρος είναι χορδή, αλλά όχι κάθε χορδή είναι διάμετρος. Καμία χορδή δεν είναι μεγαλύτερη από τη διάμετρο.
\item Το \emph{μέτρο} ενός τόξου (σε μοίρες) δεν είναι το \emph{μήκος} του. Δύο τόξα με το ίδιο μέτρο σε κύκλους με διαφορετική ακτίνα έχουν διαφορετικό μήκος.
\item Δεν υπάρχει πάντα τρίγωνο με τρία δοσμένα μήκη: οι κύκλοι μπορεί να μη τέμνονται.
\end{itemize}
\end{prosoxi}

\tmima{Λυμένα παραδείγματα}

\paradeigma{Ακτίνα και διάμετρος}
\tixero Ένας κύκλος έχει ακτίνα $3{,}5$ cm. Ένας άλλος έχει διάμετρο $5$ cm.\par
\tithelo Η διάμετρος του πρώτου και η ακτίνα του δεύτερου.\par
\possindeo $\delta=2\rho=2\cdot3{,}5=7$ cm. $\rho=\delta:2=5:2=2{,}5$ cm.\par
\apantisi $7$ cm και $2{,}5$ cm.\par
\elegxos Η διάμετρος είναι διπλάσια της ακτίνας.

\paradeigma{Μέσα, πάνω ή έξω;}
\tixero Κύκλος $(\mathrm{Ο},4)$ cm και σημεία με $\mathrm{ΟΑ}=3$ cm, $\mathrm{ΟΒ}=4$ cm, $\mathrm{ΟΓ}=5$ cm.\par
\tithelo Η θέση κάθε σημείου.\par
\possindeo Συγκρίνουμε με την ακτίνα $4$. Το $\mathrm{Α}$ έχει $3<4$, το $\mathrm{Β}$ έχει $4=4$, το $\mathrm{Γ}$ έχει $5>4$.\par
\apantisi Το $\mathrm{Α}$ είναι μέσα, το $\mathrm{Β}$ πάνω στον κύκλο, το $\mathrm{Γ}$ έξω.\par
\elegxos Μόνο το $\mathrm{Β}$ ανήκει στον κύκλο· τα $\mathrm{Α}$, $\mathrm{Β}$ ανήκουν στον κυκλικό δίσκο.

\paradeigma{Επίκεντρη γωνία και τόξο}
\tixero Σε κύκλο με κέντρο $\mathrm{Ο}$ η επίκεντρη γωνία $\widehat{\mathrm{ΑΟΒ}}$ είναι $70^\circ$.\par
\tithelo Το μέτρο του τόξου $\overset{\frown}{\mathrm{ΑΒ}}$ και του άλλου τόξου με άκρα $\mathrm{Α}$, $\mathrm{Β}$.\par
\possindeo Το τόξο που αντιστοιχεί στη γωνία έχει το ίδιο μέτρο, $70^\circ$. Το υπόλοιπο τόξο συμπληρώνει τον κύκλο των $360^\circ$.\par
\apantisi $70^\circ$ και $360^\circ-70^\circ=290^\circ$.\par
\elegxos $70^\circ+290^\circ=360^\circ$.

\paradeigma{Κατασκευή τριγώνου}
\tixero Θέλουμε τρίγωνο $\mathrm{ΑΒΓ}$ με $\mathrm{ΑΒ}=6$ cm, $\mathrm{ΑΓ}=5$ cm, $\mathrm{ΒΓ}=4$ cm. Υπάρχει τρίγωνο με πλευρές $2$, $3$ και $6$ cm;\par
\tithelo Η κατασκευή και η απάντηση για το δεύτερο τρίγωνο.\par
\possindeo Σχεδιάζουμε $\mathrm{ΑΒ}=6$ cm. Γράφουμε κύκλο $(\mathrm{Α},5)$ και κύκλο $(\mathrm{Β},4)$. Ένα από τα σημεία τομής τους είναι το $\mathrm{Γ}$. Με πλευρές $2$, $3$, $6$ θα χρειαζόμασταν τους κύκλους $(\mathrm{Α},3)$ και $(\mathrm{Β},2)$ με $\mathrm{ΑΒ}=6$. Οι δύο ακτίνες μαζί φτάνουν μόνο $3+2=5$ cm, λιγότερο από τα $6$ cm, άρα οι κύκλοι δεν τέμνονται.\par
\apantisi Το πρώτο τρίγωνο κατασκευάζεται, το δεύτερο δεν υπάρχει.\par
\elegxos Για να τέμνονται οι δύο κύκλοι πρέπει η τρίτη πλευρά να είναι μικρότερη από το άθροισμα των άλλων δύο.

\begin{center}
\begin{tikzpicture}[scale=0.55]
\draw[gray] (0,0) circle (5); \draw[gray] (6,0) circle (4);
\draw[thick] (0,0)--(6,0)--(3.75,3.31)--cycle;
\fill (0,0) circle (2pt) node[below left] {$\mathrm{Α}$};
\fill (6,0) circle (2pt) node[below right] {$\mathrm{Β}$};
\fill (3.75,3.31) circle (2pt) node[above] {$\mathrm{Γ}$};
\node at (3,-0.5) {\small $6$ cm}; \node at (1.5,2.0) {\small $5$ cm}; \node at (5.5,2.0) {\small $4$ cm};
\end{tikzpicture}
\end{center}

\paradeigma{Κύκλος που περνά από τρία σημεία}
\tixero Τρία σημεία $\mathrm{Α}$, $\mathrm{Β}$, $\mathrm{Γ}$ που δεν είναι στην ίδια ευθεία.\par
\tithelo Ο κύκλος που περνά από τα τρία.\par
\possindeo Το κέντρο $\mathrm{Ο}$ απέχει το ίδιο από το $\mathrm{Α}$ και το $\mathrm{Β}$, άρα βρίσκεται στη μεσοκάθετο του $\mathrm{ΑΒ}$. Ομοίως βρίσκεται στη μεσοκάθετο του $\mathrm{ΑΓ}$. Το κέντρο είναι το σημείο τομής των δύο μεσοκαθέτων. Ακτίνα είναι το $\mathrm{ΟΑ}$.\par
\apantisi Κέντρο η τομή των μεσοκαθέτων, ακτίνα $\mathrm{ΟΑ}$.\par
\elegxos Ο κύκλος περνά και από το $\mathrm{Γ}$, γιατί $\mathrm{ΟΓ}=\mathrm{ΟΑ}$.

\begin{center}
\begin{tikzpicture}[scale=0.8]
\draw[gblue,thick] (2,1) circle (2.236);

\draw[dashed] (2,-0.8)--(2,3.2); \draw[dashed] (-0.4,1.8)--(3.2,0.6);
\draw (0,0)--(4,0)--(1,3)--cycle;
\fill (0,0) circle (1.6pt) node[below left] {$\mathrm{Α}$};
\fill (4,0) circle (1.6pt) node[below right] {$\mathrm{Β}$};
\fill (1,3) circle (1.6pt) node[above] {$\mathrm{Γ}$};
\fill (2,1) circle (1.8pt) node[above right] {$\mathrm{Ο}$};
\end{tikzpicture}
\end{center}

\tmima{Ασκήσεις}

\zoni{Το κατάλαβα}{στοιχεία του κύκλου}

\ex Συμπλήρωσε: α) Ο κύκλος είναι ένα τόξο \ldots\ μοιρών. β) Το ημικύκλιο είναι ένα τόξο \ldots\ μοιρών. γ) Δύο κάθετες διάμετροι χωρίζουν τον κύκλο σε $4$ ίσα τόξα των \ldots\ μοιρών. δ) Αν $\rho=7$ cm, τότε $\delta=\ldots$ cm. ε) Αν $\delta=18$ cm, τότε $\rho=\ldots$ cm.

\ex Σωστό ή λάθος; α) Δύο κύκλοι είναι ίσοι όταν έχουν ίσες ακτίνες. β) Κάθε χορδή είναι διάμετρος. γ) Η διάμετρος είναι η μεγαλύτερη χορδή. δ) Όλα τα σημεία του κυκλικού δίσκου απέχουν από το κέντρο το πολύ όσο η ακτίνα. ε) Το μέτρο ενός τόξου είναι ίσο με το μέτρο της αντίστοιχης επίκεντρης γωνίας. στ) Μια χορδή μπορεί να είναι μεγαλύτερη από τη διάμετρο.

\ex Σχεδίασε ένα σημείο $\mathrm{Ο}$. Με τρία χρώματα βάψε: τα σημεία που απέχουν λιγότερο από $2$ cm από το $\mathrm{Ο}$, αυτά που απέχουν ακριβώς $2$ cm και αυτά που απέχουν περισσότερο από $2$ cm.

\ex Σχεδίασε κύκλο $(\mathrm{Ο},3\text{ cm})$ και πάρε πάνω του τρία σημεία $\mathrm{Α}$, $\mathrm{Β}$, $\mathrm{Γ}$. Σχεδίασε τη διάμετρο $\mathrm{ΑΑ'}$, την επίκεντρη γωνία $\widehat{\mathrm{ΑΟΒ}}$ και τη χορδή $\mathrm{ΒΓ}$.

\ex Δύο ομόκεντροι κύκλοι έχουν ακτίνα $2$ cm και διάμετρο $6$ cm αντίστοιχα. Βρες την ακτίνα του δεύτερου και πόσο απέχουν οι δύο κύκλοι.

\ex Ένας κύκλος έχει ακτίνα $6{,}5$ cm. Πόση είναι η διάμετρός του; Ποια είναι η μεγαλύτερη δυνατή χορδή του; Μπορεί να έχει χορδή $14$ cm;

\ex Σε κύκλο με κέντρο $\mathrm{Ο}$ είναι $\widehat{\mathrm{ΑΟΒ}}=125^\circ$. Βρες το μέτρο των δύο τόξων με άκρα $\mathrm{Α}$ και $\mathrm{Β}$.

\zoni{Το χρησιμοποιώ}{σήμα, ρολόγια, πίτσες}

\ex \textbf{Το Wi-Fi.} Ο δρομολογητής της αυλής έχει εμβέλεια $15$ m. Ο Ηλίας κάθεται $12$ m μακριά, η Ρήνη ακριβώς $15$ m και ο Στέλιος $18$ m. Ποιοι έχουν σήμα; Πόση είναι η μέγιστη απόσταση δύο ανθρώπων που έχουν σήμα όταν κάθονται αντίθετα, εκατέρωθεν του δρομολογητή;

\ex \textbf{Το ρολόι.} Η Αγγελική ζωγραφίζει ένα ρολόι. Πόσες μοίρες έχει η επίκεντρη γωνία ανάμεσα σε δύο διαδοχικούς αριθμούς; Πόσες ανάμεσα στο $12$ και στο $3$; Πόσο είναι το τόξο από το $12$ μέχρι το $5$; Πόσες μοίρες διαγράφει ο λεπτοδείκτης σε $10$ λεπτά;

\ex \textbf{Η πίτσα.} Ο Στέλιος κόβει μια στρογγυλή πίτσα σε $8$ ίσα κομμάτια. Πόση είναι η γωνία της κορυφής κάθε κομματιού στο κέντρο; Πόσο είναι το τόξο $3$ κομματιών;

\ex \textbf{Το παρτέρι.} Το κυκλικό παρτέρι του παππού Γιάννη έχει ακτίνα $2{,}4$ m. Πόσο μήκος έχει η μεγαλύτερη ευθεία διαδρομή μέσα στο παρτέρι; Μπορεί να τραβήξει σχοινί $3$ m ανάμεσα σε δύο σημεία του κύκλου; Και $5$ m;

\ex \textbf{Ο στόχος.} Ο Βασίλης παίζει βελάκια σε στόχο με ομόκεντρους κύκλους ακτίνας $5$, $10$ και $15$ cm. Το βελάκι του καρφώνεται $12$ cm από το κέντρο. Ανάμεσα σε ποιους δύο κύκλους βρίσκεται;

\zoni{Το σκέφτομαι}{κατασκευές}

\ex α) Κατασκεύασε τρίγωνο με πλευρές $5$, $4$ και $3$ cm. β) Υπάρχει τρίγωνο με πλευρές $7$, $2$ και $4$ cm; Δικαιολόγησε.

\ex Κατασκεύασε ισόπλευρο τρίγωνο πλευράς $5$ cm μόνο με χάρακα και διαβήτη. Εξήγησε γιατί οι τρεις πλευρές του είναι ίσες.

\ex Η Ζωή βρήκε ένα σπασμένο κυκλικό πιάτο και θέλει να βρει το κέντρο του. Πώς μπορεί, με τρία σημεία στην άκρη του πιάτου και διαβήτη, να το βρει;

\ex Η Πολυξένη σχεδιάζει κύκλο ακτίνας $3$ cm και, με το διαβήτη ανοιχτό $3$ cm, πηγαίνει από σημείο σε σημείο πάνω στον κύκλο. Μετά από πόσα βήματα επιστρέφει στην αρχή; Εξήγησε με τα τρίγωνα που σχηματίζονται.

\zoni{Για το σπίτι}{ρίξε μια ματιά}

\ex Βρες τρία κυκλικά αντικείμενα στο σπίτι. Μέτρησε τη διάμετρό τους και βρες την ακτίνα τους. Πώς θα βρεις το κέντρο ενός χάρτινου κύκλου διπλώνοντάς τον;

\ex Σχεδίασε κύκλο με διάμετρο $7$ cm, χώρισέ τον σε $6$ ίσα τόξα και φτιάξε ένα διακοσμητικό «λουλούδι».
